

Par
ailleurs, $j$ définit une section $j'$ de $X'=V'\times_V X$:
\[
\begin{tikzcd}[column sep=large,row sep=large]
X' \arrow[r,"i'"] \arrow[d,"p'"'] & V' \arrow[d,"p"]\\
X \arrow[r,"i"] \arrow[ur,"j"] \arrow[u,"j'"',bend right] & V .
\end{tikzcd}
\]
Soient $\underline{I},\underline{J},\underline{K}$ les idéaux définissant $i,j,j'$ ; on a : $j=i'\circ j'$, et l'immersion $j'$ est régulière; de plus, $i'$ est définie par l'idéal $p^*(\underline{I})$ (car $p$ est plat), et son faisceau conormal est donc $i'^*(p^*(\underline{I}))=p'^*(\underline{I}/\underline{I}^2)$. On obtient alors la suite exacte :
\begin{equation}\tag{2.2.3}
0\longrightarrow \underline{I}/\underline{I}^2\longrightarrow \underline{J}/\underline{J}^2\longrightarrow \underline{K}/\underline{K}^2\longrightarrow0.
\end{equation}
Nous allons montrer :

a) que le diagramme (2.2.1) est commutatif ; on obtiendra donc ainsi un morphisme de complexes $\theta$ du complexe $\underline{L}^{X/Y}_\bullet$ défini par $V$ dans le complexe $\underline{L}^{X/Y}_\bullet$ défini par $V'$ ;

b) que l'homomorphisme : $\underline{K}/\underline{K}^2\longrightarrow j^*(\underline{\Omega}^1_{V'/V})$ défini par (2.2.2), (2.2.3) et (2.2.1) est un isomorphisme ; cela entraînera que le mapping cylinder de $\theta$ est acyclique, donc que $\theta$ est un quasi-isomorphisme, donc définit un isomorphisme dans la catégorie dérivée.

Les propriétés a) et b) sont de nature locale sur $X$. On peut donc supposer $X,Y,V$ et $V'$ affines. Posons : $Y=\Spec(A)$, $V=\Spec(B)$, $V'=\Spec(B')$, $X=\Spec(C)$, et gardons les mêmes lettres pour noter les homomorphismes et les idéaux correspondants. On a le diagramme :
